\clearpage
\section{Overview}\label{page:overview}

\subsection{Contract Position}
\BedrockContractDiagram{foundation}

\subsection{Manual Scope}
This manual defines programmer-visible architectural state, operand encoding, execution semantics,
and instruction behavior for the Bedrock instruction set architecture.

\Needspace{3.5in}
\subsection{Architectural Profile}
Bedrock defines sixteen 64-bit general-purpose registers, R0 through R15. SP and PC are special architectural registers outside the
Rn namespace, with SP-relative addressing tied to SS and PC-relative addressing tied to CS.

\begin{itemize}
\item The PC names a byte in the instruction stream; sequential execution advances it by the encoded instruction length determined during decoding.
\item Encoded instruction length is explicit and bounded; the base profile defines encoded instruction lengths from 1 to 18 bytes.
\item Four integer operation sizes: B, W, L, and Q.
\item The optional scalable-vector extension provides 32 VLEN-bit vector registers and 16 packed predicate registers; VLEN is a reset-stable power of two from 128 through 2048 bits.
\item The effective-address model covers register, memory, immediate, absolute, segment-qualified, indexed, and auto-update operands.
\item Segment pre-translation precedes optional page-table translation.
\item User and supervisor modes share the instruction set, with privileged instructions and checked user-access moves defining supervisor-only behavior.
\end{itemize}
